## TEMP holding note — RQ3 outlier diagnosis, the actual point of building SHAP (2026-08-12)

**New**: `models/growth_curve_attribution/rq3_outlier_diagnosis.py` — pulls the N largest-residual
plots and shows, per plot, which features' SHAP contributions drove that specific prediction.
Purely local, reads only already-saved `predictions.csv` + `xgboost_shap_values.csv`, no fitting.
Run for real: `--cohort 4survey --set-name nested_set4_gated_all_vif --top-n 5`.

### Results — top 5 worst XGBoost residuals, Set4/4survey

| identification | cpmt | observed | predicted | residual | Top SHAP driver |
|---|---|---:|---:|---:|---|
| 353952 | 2232 | 59.12 | 2.78 | +56.34 | `CanopyCover` -3.95, `tas_mean` +1.32, `chelsa_bio12_precip_mm` +1.24 |
| 77185 | 1135 | -52.71 | 3.36 | -56.07 | `CanopyCover` -6.08, `chelsa_bio12_precip_mm` +2.87 |
| 356411 | 2217 | 60.66 | 6.00 | +54.66 | `chelsa_bio12_precip_mm` +3.41, `cpmt_compactness_ratio` -1.27 |
| 362556 | 2173 | 56.87 | 3.80 | +53.08 | `chelsa_bio12_precip_mm` +0.82, `slope_degrees` +0.56 |
| 370819 | 2176 | 58.20 | 5.70 | +52.50 | `chelsa_bio12_precip_mm` +1.96, `topex` +0.49 |

### Extended to all 6 (set x cohort) combos, 2026-08-13

Same script, `--top-n 5`, run for the remaining 5 combos (Set2/Set3 x both cohorts). **Result: the
same handful of plots recur across nearly every combo, not 6 independent sets of outliers.**
`identification` 356411, 353952, and 77185 appear in the worst-5 for literally all 6 combos
(3 sets x both cohorts they belong to); 344808 and 370819 recur in most of the 4survey combos.
This is itself informative: the outlier-ness is a property of the PLOT, not of which feature set
or model happens to be fit — ruling out "one set's specific variables are causing spurious
predictions" as an explanation.

### Cross-reference against the disturbance-classification system (2026-08-13) — CORRECTS the
### "almost certainly disturbance/data-quality artifacts" claim below

New: `models/growth_curve_attribution/cross_reference_outliers_with_disturbance.py` — joins the
10 unique outlier `identification`s (across all 6 combos) against
`disturbance_checks.summarize_plot_disturbance_status()` (the same `any_clearfell_like`/
`any_measurement_inconsistent`/`any_ambiguous_disturbance` flags the notebook's own compartment
map uses) and `fit_all_years_and_compute_residuals()`'s `residual_range` (the same quantity behind
the notebook's `extreme_trajectory_flag`, top-1% cutoff). Calls the standalone functions in
`disturbance_checks.py` directly — the notebook itself was never opened or executed.

**Result: 0 of 10 outlier plots carry ANY of the three disturbance/measurement flags, and 0 of 10
clear the top-1% `extreme_trajectory_flag` cutoff either** (residual_range for all 10 plots is
0.9-15.1m, cohort cutoffs are 17.8m/4survey and 23.2m/6survey — every one sits below its own
cohort's threshold). **The "almost certainly disturbance/data-quality artifacts" framing below was
not actually checked against real data when first written, and turns out to be wrong** — this is
exactly the kind of claim that needed verifying before being cited, and it doesn't hold up.

**A more specific, better-supported pattern emerges from the residual_range split itself**:
- **5 of 10 plots** (344808, 353952, 356411, 344812, 362556) have moderately-high residual_range
  (13.6-15.1m) — real within-trajectory curve-fit instability across their own survey years, just
  not extreme enough to clear the top-1% bar or match the specific consecutive-survey collapse
  patterns the classifier looks for.
- **5 of 10 plots** (77185, 370819, 320299, 366554, 366559) have LOW residual_range (0.9-5.9m) —
  meaning these plots' own height trajectory fits a SMOOTH, internally consistent curve across
  every survey year. Their entire `local_y_max_difference` (40-60m) comes from that smooth curve
  sitting far away from the *official yldc benchmark*, not from any instability in the
  LiDAR-observed growth itself.

**Revised interpretation**: for the low-residual_range subset, the most defensible explanation is
no longer "disturbance" but a **yldc lookup mismatch** — e.g. the wrong yield class or species
assumed for that stand in the official table, which would produce exactly this signature (smooth,
internally consistent observed growth; large, stable offset from the benchmark curve). For the
higher-residual_range subset, "real but sub-threshold trajectory instability" is a more honest
label than "disturbance," since the existing classifier explicitly does not flag them. Neither
group supports the original blanket "data-quality artifact" claim — extreme `local_y_max_difference`
values on these specific plots look like genuine signal about the yldc benchmark or the site
itself, not a case to exclude from RQ3's attribution numbers.

**Practical implication for RQ3's headline numbers**: no principled basis for excluding these
plots was found — they should stay in the modelling population as-is. The earlier framing that
implied excluding them "inflates the unexplained residual story" no longer applies; there's no
verified evidence these are the artifacts they were assumed to be.

### Not yet done

Checking whether these same identification numbers are also outliers in EN/GNNWR predictions
(would test whether the yldc-mismatch/sub-threshold-instability read generalizes beyond XGBoost);
directly inspecting the yldc lookup (yield class/species assumptions) for the 4 low-residual_range
plots to test the yldc-mismatch hypothesis concretely, rather than leaving it as inference from the
residual pattern alone.
